% Contents/status_e_condicoes.tex
\midsectiontitle{Estados e Condições}

\begin{rpgboxtips}
    {Nota para o Aventureiro} \textbf{Fique Atento:} Se você não lembra dos
    status aplicados do seu personagem, o Mestre não tem a obrigação de lembrá-los
    por você. Condições são provenientes de magias, ataques, poções e elixires, e
    podem exigir rolamentos específicos dependendo da situação.
\end{rpgboxtips}

\titlesection{Status Negativos}

\originlink{status:abalado}
\pericia{Abalado}{O espírito do alvo está fraquejado. Todos os testes realizados pelo personagem sofrem \textbf{Desvantagem}. O alvo deve realizar um teste de \link{att:sense}{Sense} para quebrar o efeito.}

\pericia{Atordoado}{O alvo fica incapacitado de agir e reagir. Após sofrer este efeito, o alvo torna-se imune a novos atordoamentos de mesma origem até o fim da cena. Caso tenha resistido ao atordoamento, recebe \textbf{1D4} rodadas de imunidade contra esta condição.}

\pericia{Aterrorizado}{O alvo vai gastar todas as suas ações para fugir da fonte do medo por \textbf{1D4} rodadas. Exige um teste de \link{att:sense}{Sense} para quebrar o efeito.}

\pericia{Cegueira}{A escuridão toma os olhos. O alvo recebe \textbf{Desvantagem} em todos os testes de Ataque, Magia e Percepção.}

\originlink{status:concussionado} 
\pericia{Concussão}{O impacto desorientou os sentidos. O alvo deve escolher entre manter sua \textbf{Ação de Movimento} ou sua \textbf{Reação}. Exige um teste de \link{att:fortitude}{Fortitude} contra o valor do ataque sofrido.}

\pericia{Congelado / Petrificado}{O corpo está rígido como gelo ou rocha. O alvo deve gastar sua \textbf{Ação de Movimento} para se libertar, ou ficará impossibilitado de esquivar, apenas bloquear.}

\pericia{Controlado}{A vontade do atacante sobrepuja a do alvo. O personagem agirá conforme os comandos de quem o controla. Exige um teste de \link{att:sense}{Sense} para quebrar o encantamento.}

\pericia{Debilitado}{Causado ao sofrer um dano massivo (\textbf{90\% da Vida Máxima} em um único golpe). O alvo deve escolher entre realizar uma Ação Padrão ou uma Ação de Movimento no seu turno.}
\originlink{status:derrubado} \pericia{Derrubado}{O alvo foi ao chão. Deve gastar sua \textbf{Ação de Movimento} para se levantar.}

\originlink{status:desprevenido}
\pericia{Desprevenido}{A guarda está baixa. Alvos nesta condição permitem que o primeiro ataque recebido utilize regras de \textit{Furtividade} (Ataque Furtivo).}

\pericia{Doente}{O personagem sofre de enfermidades mágicas ou biológicas, sofrendo penalidades constantes definidas pela doença.}

\pericia{Encantado}{O alvo obedece à pessoa que causou o efeito. Exige um teste de \link{att:sense}{Sense} para resistir inicialmente e para quebrar o efeito nas rodadas seguintes.}

\pericia{Exaurido}{O oposto de Revigorado. O alvo perde o dobro de IEP ao usar habilidades e não recupera IEP passivamente.}

\originlink{status:encharcado} \pericia{Encharcado}{O alvo está coberto por água. Recebe \textbf{Resistência (Dano pela metade)} contra Fogo, mas sofre \textbf{Dano Multiplicado (x1.5)} contra Raio e Eletricidade.}

\pericia{Envenenado}{O alvo sofre os efeitos do veneno a cada início de turno. Deve realizar um teste de \link{att:fortitude}{Fortitude} contra a toxicidade do veneno em cada rodada, ou consumir antídotos para tratar.}

\pericia{Flanqueado}{Cercado por inimigos. Reduz a Defesa em \textbf{5 pontos} e o alvo é considerado \textbf{Desprevenido}.}
\pericia{Lento}{Alvos lentos não podem saltar ou voar. Sua \textbf{Ação de Movimento} é reduzida pela metade.}
\pericia{Inconsciente}{O alvo está indefeso e totalmente incapaz de agir ou reagir.}

\pericia{Perplexo}{O choque mental impede ações ofensivas. O alvo não pode agir, mas consegue se defender normalmente.}
\pericia{Vulnerável}{O alvo perde sua Redução de Dano (Resistência) por uma rodada.}

\pericia{Provocado}{A fúria cega o julgamento. O alvo deve atacar quem o provocou até a morte. Exige um teste de \link{att:sense}{Sense} + \link{per:temperanca}{Temperança} para resistir.}
\originlink{status:queimado} \pericia{Queimado}{O alvo está em chamas. Deve gastar sua \textbf{Ação de Movimento} para apagar o fogo ou sofrerá dano contínuo até o efeito ser cessado.}
\originlink{status:ardiloso} \pericia{Sangramento Ardiloso}{Dano contínuo por acúmulo. Para cada \textit{stack}, o alvo sofre 1D4. A remoção de cargas depende da sua robustez física:}

\begin{center}
    \scriptsize
    \renewcommand{\arraystretch}{1.3}
    \begin{tabularx}
        {\columnwidth}{|X|c|c|c|c|c|} \hline \rowcolor{overgrownRed!10} \multicolumn{6}{|c|}{\textbf{Escalonamento
        por Grace/Wisdom}} \\
        \hline \textbf{Gra/Wis} & 10 & 20 & 40 & 80 & 100 \\
        \hline \textbf{Dano} & $1 \times D4$ & $1 \times D8$ & $1 \times D10$ & $1
        \times 12$ & $1 \times D20$ \\
        \hline
    \end{tabularx}
\end{center}

\begin{center}
    \scriptsize
    \renewcommand{\arraystretch}{1.3}
    \setlength{\tabcolsep}{3pt}
    \begin{tabularx}
        {\columnwidth}{l @{\extracolsep{\fill}} ccccccc} \toprule \rowcolor{gray!20}
        \textbf{Fortitude} & 15 & 30 & 45 & 60 & 75 & 90 & 100 \\
        \midrule \textbf{Stacks Removidos} & 2 & 4 & 6 & 8 & 10 & 12 & 14 \\
        \bottomrule
    \end{tabularx}
\end{center}
\originlink{status:brutal} \pericia{Sangramento Brutal}{Um ferimento profundo e único. Em ataques bem-sucedidos, causa $2D8 + \text{Might/Wisdom}- \text{Fortitude}$. O dano bruto escala conforme a maestria do atacante:}

\begin{center}
    \scriptsize
    \renewcommand{\arraystretch}{1.4} % Aumenta levemente a altura das linhas
    \begin{tabularx}
        {\columnwidth}{|l|X|X|X|X|} \hline \rowcolor{overgrownRed!10} \multicolumn{5}{|c|}{\textbf{Escalonamento
        por Might/Wisdom}} \\
        \hline \textbf{Might/Wisdom} &
        \centering
        25 &
        \centering
        50 &
        \centering
        75 &
        \centering
        100 \tabularnewline \hline \textbf{Dano} &
        \centering
        $2 \times D8$ &
        \centering
        $3 \times D10$ &
        \centering
        $4 \times D10$ &
        \centering
        $8 \times D12$ \tabularnewline \hline
    \end{tabularx}
\end{center}
\originlink{status:segurado} \pericia{Segurado / Preso}{O movimento é impossibilitado. O alvo deve realizar um teste para se libertar.}

\originlink{status:silenciado}
\pericia{Silenciado}{O fluxo de energia foi interrompido. O alvo não pode conjurar magias, mas ainda pode utilizar itens mágicos. Exige teste de \link{att:wisdom}{Wisdom} para parar.}

\pericia{Sobrecarregado}{O alvo sobrecarregado tem desvantagem em testes de luta e movimentação reduzida.}
\originlink{status:sufocado} \pericia{Sufocado}{O alvo sufocado tem desvantagem em testes de luta, esquiva e movimentação reduzida.}

\titlesection{Status Positivos}

As vantagens abaixo são provenientes de fontes de \textit{buff} (magias, itens
ou habilidades).

\pericia{Aprimorado}{Sentidos aguçados ao extremo. O alvo tem \textbf{Vantagem} em testes de \link{att:sense}{Sense} até o fim da cena.}

\pericia{Celeridade}{O corpo move-se além dos limites. O alvo tem \textbf{Vantagem} em testes de \link{att:grace}{Grace} até o fim da cena.}

\pericia{Esclarecido}{Compreensão absoluta do ambiente. O alvo tem \textbf{Vantagem} em testes de \link{att:wisdom}{Wisdom} até o fim da cena.}

\pericia{Fortalecido}{Músculos levados ao limite superior. O alvo tem \textbf{Vantagem} em testes de \link{att:might}{Might} até o fim da cena.}

\pericia{Inspirado}{O alvo recebe um bônus de \textbf{1D8} em uma rolagem à sua escolha ou pode adicionar um dado extra de dano.}

\pericia{Invisível}{Presença indetectável. Todos os ataques são considerados furtivos e atingem alvos \textbf{Desprevenidos}.}

\pericia{Passo Leve}{Movimentação silenciosa e fluida. O alvo tem \textbf{Vantagem} em testes de \link{per:furtividade}{Furtividade}.}

\originlink{status:respiracaosubmersa}
\pericia{Respiração Submersa}{Capacidade mágica ou biológica de respirar debaixo d'água.}

\pericia{Revigorado}{Capacidade física aumentada. O alvo tem \textbf{Vantagem} em testes de \link{att:fortitude}{Fortitude}.}

\pericia{Acelerado}{Ganha 1 ação padrão.}